\documentclass[compacto,10pt]{pubciencias-notas}

% -- Paquetes adicionales
\usepackage{comandosEPN}

% -- Datos del libro
\autor{Departamento de Formación Básica}
\materia{Álgebra Lineal}
\nota{Semana 10 Clase 02}
\fecha{Semestre 2025-A}

\logouno[1.1cm]{Logo_EPN.eps}
\logodos[1.7cm]{Logo_DFB.eps}

\definecolor{colordef}{cmyk}{0.81,0.62,0.00,0.22}

\begin{document}

\encabezado

\vspace*{-4mm}

\setcounter{section}{3}
\section{Espacios Vectoriales}
\setcounter{subsection}{7}
\subsection{Suma de espacios vectoriales}

\begin{teo}
    Sean $(E,+,\cdot,\K)$ un espacio vectorial de dimensión finita. Para todo $S$ subespacio vectorial de $E$, existe un subespacio vectorial $T$ de $E$ tal que 
    \[
        E = S \oplus T.
    \]
    Se dice que $T$ es el subespacio complementario de $S$ en $E$.
\end{teo}
\begin{ejem}\hspace{0pt}\label{EjemDirecta2} En el espacio vectorial $\R^{3}$ considere los subespacios vectoriales
\[
    W_{1} = \{(x,y,0)\in \R^3: x,y \in \R \}
    \texty 
    W_{2} = \{(0,0,z)\in \R^3: z \in \R\}.
\]
Se tiene que $W_2$, es el subespacio vectorial complementario de $W_1$ en $\R^{3}$.
\end{ejem}

\begin{ejer}\label{EjercicioSumaDirecta1}\hspace{0pt} Sea $W=\left\{(a,b,c,d)\in \mathbb{R}^{4} : a+b+c=0 \right\}$ un subespacio vectorial de $\mathbb{R}^{4}$. Encuentre un subespacio vectorial $U$ tal que:
\[\mathbb{R}^{4} = W \oplus U\]
\end{ejer}


\begin{teo}\label{TeoremaSumaDimensión}
Sean $(E,+,\cdot,\K)$ un espacio vectorial de dimensión finita y $S,T\subseteq E$ subespacios vectoriales de $E$. Si $E$ es la suma directa de $S$ y $T$, entonces
\[
    \dim(E) = \dim(S) + \dim(T).
\]
\end{teo}

\begin{ejem}\hspace{0pt} Considere el ejemplo \ref{EjemDirecta2}, de acuerdo con el teorema anterior se verifica que:
\[
    \dim(\R^{3}) = \dim(W_{1}) + \dim(W_{2}).
\]
En efecto, se puede verificar que:
\[
    \dim(W_{1}) = 2
    \texty
    \dim(W_{2}) = 1.
\]
Notemos además que $W_{1} \cap W_{2} = \{0_{\R^{3}}\}$.
\end{ejem}

\begin{ejer}\hspace{0pt} Verifique que se satisface el Teorema~\ref{TeoremaSumaDimensión} para el Ejercicio~\ref{EjercicioSumaDirecta1}.
    
\end{ejer}

\end{document}